\item \textbf{¿Por qué no es posible la siguiente situación?}
Para estudiar el positronio, un científico coloca muestras sólidas de $^{57}\text{Co}$ y $^{14}\text{C}$ en estrecha proximidad física. El $^{57}\text{Co}$ decae por emisión de $e^+$ y el $^{14}\text{C}$ decae emitiendo $e^-$. Sostiene que la combinación directa de los positrones y electrones emitidos produce una cantidad abundante de átomos de positronio estables para su investigación.